\begin{problem}[第35届全国中学生物理竞赛复赛][塞曼效应与洛伦兹解释]{92}
五、塞曼发现了钠光D线在磁场中分裂成三条，洛伦兹根据经典电磁理论对此做出了解释，他们因此荣获1902年诺贝尔物理学奖。假定原子中的价电子（质量为 $m$，电荷量为 $-e$，$e>0$）受到一指向原子中心的等效线性回复力 $-m\omega_{0}^{2}r$ （$r$ 为价电子相对于原子中心的位矢）作用，做固有圆频率为 $\omega_{0}$ 的简谐振动，发出圆频率为 $\omega_{0}$ 的光。现将该原子置于沿 $z$ 轴正方向的匀强磁场中，磁感应强度大小为 $B$（为方便起见，将 $B$ 参数化为 $B=\frac{2m}{e}\omega_{L}$ ）。
\begin{subproblem}
选一绕磁场方向匀角速转动的参考系，使价电子在该参考系中做简谐振动，导出该电子运动的动力学方程在直角坐标系中的分量形式并求出其解；
\end{subproblem}
\begin{subproblem}
将（1）问中解在直角坐标系中的分量形式变换至实验室参考系的直角坐标系；
\end{subproblem}
\begin{subproblem}
证明在实验室参考系中原子发出的圆频率为 $\omega_{0}$ 的谱线在磁场中一分为三；并对弱磁场（即 $\omega_{L} \ll \omega_{0}$ ）情形，求出三条谱线的频率间隔。

已知：在转动角速度为 $\omega$ 的转动参考系中，运动电子受到的惯性力除惯性离心力外还受到科里奥利力作用，当电子相对于转动参考系运动速度为 $v'$ 时，作用于电子的科里奥利力为 $f_{c} = -2m\omega \times v'$ 。
\end{subproblem}
\end{problem}

